\begin{table}[!tbp]
\centering
\small
\caption{Contraste con cifras oficiales y verificación de fuga}
\label{tab:a4}
\begin{tabular}{p{4.6cm}p{2.7cm}p{5.0cm}p{1.2cm}}
\toprule
Concepto & Archivo & Referencia & Dif. \\
\midrule
\addlinespace[2pt]\multicolumn{4}{l}{\textit{Contraste con cifras oficiales, enero a marzo de 2026}}\\
Siniestros & 4.789 & INEC: 4.789 & 0\,\% \\
Lesionados & 4.032 & INEC: 4.032 & 0\,\% \\
Fallecidos & 603 & INEC: 603 & 0\,\% \\
\addlinespace
Siniestros, enero a mayo de 2026 & 8.718 & ANT: 8.718 & 0\,\% \\
\addlinespace[2pt]\multicolumn{4}{l}{\textit{Verificación de fuga de información}}\\
Condición \texttt{FALLECIDO} y columna FALLECIDOS & 17.236 de 17.236 & Coincidencia exacta & 100\,\% \\
Condición \texttt{LESIONADO} y columna LESIONADOS & 17.236 de 17.236 & Coincidencia exacta & 100\,\% \\
Categorías de la columna CONDICIÓN & 4 & FALLECIDO, ILESO, LESIONADO, NO IDENTIFICADO & --- \\
\bottomrule
\end{tabular}
\par\vspace{2pt}{\footnotesize Referencias: INEC [1] y ANT [5] para los totales; [2] para el archivo. La segunda parte muestra que la condición de los participantes contiene el desenlace y no puede usarse como entrada. Cálculo propio (Anexo B).}
\end{table}
